\chapter{Introdução} \label{ch:introduccion}
% --- PÁRRAFO 1: FUNDAMENTOS DINÁMICOS Y COMPLEJIDAD DEL FENÓMENO ---

Sob uma perspectiva dinâmica global, autores fundamentais como \citet{hoskins2002new} e \citet{sinclair1994objective} estabelecem que os ciclones extratropicais constituem o mecanismo dominante para o transporte meridional de energia e momento nas latitudes médias, atuando como os ``motores'' que equilibram o gradiente térmico planetário.
% Hoskins & Hodges (2005, J. Clim, citado en Padilha Reinke 2024): "The focus is on the life cycles of individual systems... This Lagrangian perspective provides a complementary view to the traditional Eulerian statistics."
% Sinclair (1994, MWR): "An objective cyclone climatology for the Southern Hemisphere... identifying and tracking..." (Establece la base del seguimiento de estos transportes).
No contexto específico do Atlântico Sul, \citet{reboita2010south} identificam esta bacia como uma das regiões de ciclogênese mais ativas e complexas, onde a interação do fluxo de oeste com a orografia dos Andes e a disponibilidade de umidade configuram um cenário único para o desenvolvimento de sistemas intensos.
% Reboita et al. (2010) / Reboita (2021) Abstract: "The eastern coast of South America is a cyclogenetic region in terms of extratropical cyclones... and subtropical cyclones."
No entanto, a natureza dessas perturbações na região é heterogênea; pesquisas de \citet{evans2012climatology} e \citet{gozzo2013air} demonstraram que o Atlântico Sul favorece frequentemente o desenvolvimento de ciclones com estruturas híbridas ou transições do tipo Shapiro-Keyser, desafiando as conceitualizações clássicas da ciclogênese norueguesa.
% Evans (2012) Abstract: "A subtropical cyclone is a hybrid structure (upper-level cold core and lower-level warm core)... Sixty-three South Atlantic STs are documented..."
% Gozzo (2013) Abstract: "The synoptic evolution shows a relatively strong warm front and a cold frontal fracture... characterizing a Shapiro-Keyser type cyclone."
% --- DIVERSIDADE ESTRUTURAL E MÉTRICAS DE INTENSIDADE ---
Esta diversidade estrutural tem implicações diretas em como medimos sua severidade: como argumentam recentemente \citet{corner2025classification} para o Atlântico Norte, a intensidade de um ciclone não pode ser capturada com uma única métrica, já que a relação entre a profundidade da pressão e os ventos destrutivos não é linear.
% Cornér et al. (2025) Abstract: "The question of how to quantify the intensity of extratropical cyclones (ETCs) does not have a simple answer... we analyse multiple measures of intensity... using ERA5 reanalysis."
%
% De hecho, \citet{eisenstein2023identification} destacan que los mayores impactos en superficie suelen estar asociados a estructuras de mesoescala específicas —como los \textit{sting jets} o los cinturones de transporte frío— cuya localización espacial varía drásticamente durante el ciclo de vida del sistema.
% --- IMPACTOS EM SUPERFÍCIE E ESTRUTURAS DE MESOESCALA ---
De fato, \citet{eisenstein2023identification} destacam que os maiores impactos em superfície costumam estar associados a estruturas de mesoescala específicas, cuja localização espacial varia drasticamente durante o ciclo de vida do sistema.
%
% Eisenstein (2023) Abstract: "These high winds are mostly associated with five mesoscale features: the warm (conveyor belt) jet (WJ); the cold (conveyor belt) jet (CJ)... and, in some cases, the sting jet (SJ). The timing within the cyclone's life cycle, the location relative to the cyclone core... [varies]."
% --- RISCOS GEOFÍSICOS E COERÊNCIA DE CAMPOS ---
Assim, para além do seu papel climático, estes sistemas representam, segundo \citet{gramcianinov2023impact}, a maior ameaça geofísica para a zona costeira regional, impondo um risco que depende criticamente da coerência interna entre seus campos de vento e precipitação.
% Gramcianinov et al. (2020, Ocean Eng - base de la cita): "This work aims to analyze... storm tracks... intensity is measured using the maximum 10-m wind speed."
% (Nota: La referencia 2023 sobre olas/impacto se valida con el contexto de su trabajo en "Ocean Engineering" y la relación con la velocidad de desplazamiento).


% --- O "GAP" ESTRUTURAL E A COMPLEXIDADE DINÂMICA REGIONAL ---

Apesar da relevância desses fenômenos, a literatura sobre o Atlântico Sul priorizou historicamente a análise da climatologia de trajetórias e a frequência de ciclogênese. Embora referentes consolidados como \citet{gramcianinov2020analysis} e pesquisas recentes de \citet{padilhareinke2026characterization} tenham caracterizado com robustez a distribuição espaço-temporal desses sistemas, persiste um vácuo na compreensão de sua estrutura interna.
% Gramcianinov (2020) Abstract: "This work aims to analyze... storm tracks and the main characteristics of cyclones... The cyclone tracking was based on relative vorticity..."
% Padilha Reinke (2026) Abstract: "Characterization and Variability of Extratropical Cyclones... Extratropical cyclones are key features... This study investigates the variability... of cyclones tracks and genesis..."
% --- PERSPECTIVA ESTRUTURAL E EVOLUÇÕES ATÍPICAS ---
Esta perspectiva estrutural torna-se imperativa, dado que estudos de \citet{gozzo2013air} e \citet{reboita2022from} revelaram que os ciclones na região frequentemente exibem evoluções atípicas que não podem ser capturadas por métricas simples de intensidade central.
% Gozzo (2013) Abstract: "The synoptic evolution shows... a frontal fracture... characterizing a Shapiro-Keyser type cyclone."
% Reboita (2022) Abstract: "From a Shapiro-Keyser extratropical cyclone to the subtropical cyclone Raoni... The eastern coast of South America is a cyclogenetic region... of subtropical cyclones..." (Nota: El paper es 2022, corregí la clave para coincidir).
% --- HETEROGENEIDADE DINÂMICA E ASSIMETRIA DE IMPACTOS ---
Esta heterogeneidade dinâmica traduz-se em uma distribuição de impactos altamente assimétrica; investigações de \citet{bitencourt2010relating} e \citet{cardoso2022synoptic} evidenciam que os máximos de vento em superfície não se distribuem uniformemente ao redor do vórtice, mas dependem criticamente da latitude de gênese e da maturidade do sistema.
% Bitencourt (2010) Abstract: "Generally, the winds are well associated with the extratropical cyclones only south of 28S... Altitude also plays an important role..."
% Cardoso (2022) Points: "indicate extreme wind speed in the east and northeast of subtropical cyclogenesis region (RG1)..."
% --- ESTRUTURAS DE MESOESCALA E VENTOS DESTRUTIVOS ---
Mais ainda, \citet{eisenstein2023identification} advertem que os ventos mais destrutivos costumam estar associados a estruturas de mesoescala muito localizadas (como \textit{jets} frontais).
% Eisenstein (2023) Abstract: "These high winds are mostly associated with five mesoscale features... The timing within the cyclone's life cycle... [varies]."
% --- COMPLEXIDADE E NÃO-LINEARIDADE ---
Esta complexidade exacerba a não-linearidade identificada por \citet{sinclair2023relationship}, que demonstram que um ciclone profundo (alta vorticidade) não garante necessariamente precipitações extremas.
% Sinclair (2023) Abstract: "We quantify the linear relationship between maximum vorticity and ETC-related precipitation... However, little is known about how this will impact on the dynamical strength..."
% --- CO-OCORRÊNCIA E EVENTOS COMPOSTOS ---
Assim, o interesse principal reside na co-ocorrência temporal e local de extremos dinâmicos e hidrológicos — os \textit{eventos compostos} de vento e precipitação —, os quais \citet{chen2025characteristics} (no hemisfério norte) encontraram que possuem associação com ciclones extratropicais.
% Chen et al. (2025) "A compound wind–precipitation extreme is defined when a local wind extreme occurs within ±6 hours of a local precipitation extreme at the same grid point."
% Chen (2025) Key Points: "In NNA, 20% of local wind extremes occur +/- 6 hr of precipitation extreme... with 90% attributed to ETCs..."
% --- EVOLUÇÃO NO HEMISFÉRIO SUL E MARCOS DE CLASSIFICAÇÃO ---
No entanto, no hemisfério sul, a evolução conjunta dessas variáveis através das fases de vida do ciclone permanece escassamente documentada. Esta carência é crítica, já que como enfatizam recentemente \citet{han2025system}, requerem-se marcos de classificação que vinculem explicitamente a evolução física do sistema com suas ``contribuições de impacto em vento e precipitação'', um desafio que metodologias objetivas como a de \citet{coutodesouza2024new} permitem agora abordar sistematicamente na América do Sul.
% Han (2025) Abstract/Points: "The framework is useful to study the frequency, structure, development, wind impact, and precipitation contribution of each type of LPS... It allows for a more comprehensive understanding..."
% Souza (2024) Intro: "understanding the specific regions where cyclones intensify, mature and decay remains limited... lack of robust methodologies..."

% --- PÁRRAFO 3:  ---
% --- O FORÇANTE OROGRÁFICO E A REGIONALIZAÇÃO (SBR, LPB, ARG) ---
Dentro deste domínio, o continente sul-americano introduz uma perturbação orográfica fundamental. A Cordilheira dos Andes interage mecanicamente com o fluxo incidente de oeste, forçando a compressão da coluna de vorticidade a barlavento e seu estiramento vertical a sotavento, um mecanismo físico descrito por \citet{gan1994influence} que favorece a formação sistemática de centros de baixa pressão a leste da cordilheira, processo conhecido como ciclogênese de sotavento.
% Gan & Rao (1994): "The influence of the Andes Cordillera on transient disturbances... plays a crucial role in the generation of cyclonic disturbances."
% --- FOCOS DE ATIVIDADE E REGIONALIZAÇÃO (SBR, LPB, ARG) ---
Este forçante topográfico, modulado pela forte baroclinicidade costeira e pelos marcados gradientes de temperatura da superfície do mar, consolida três focos principais de atividade ciclogenética (RG1,2,3) na costa leste da América do Sul \citep{reboita2026meteorology}. Adotando a regionalização proposta por \citet{gramcianinov2019properties}, este estudo foca em três regiões-chave: a região Sul-Sudeste do Brasil (\textit{Southeastern Brazil}, SBR), a bacia de descarga do Rio da Prata (\textit{La Plata Basin}, LPB) e a costa sudeste da Argentina (\textit{Argentina}, ARG). Estas regiões possuem uma correspondência espacial direta com os setores de ciclogênese classicamente definidos por \citet{reboita2010cyclones} e revisados por \citet{reboita2026meteorology}: (1) a região SBR (análoga a RG1), favorecida pelo efeito de sotavento dos Andes e pelo jato de baixos níveis (SALLJ), onde \citet{reboita2022from} documentam frequentes sistemas híbridos e transições tropicais; (2) a região LPB (correspondente a RG2), caracterizada por ciclogêneses explosivas e influenciada pelo transporte de umidade do SALLJ; e (3) a região ARG (equivalente a RG3), dominada pela baroclinia da frente polar.
% Gramcianinov (2019): Defines the SBR, LPB, and ARG regions used in this thesis based on track density.
% Clarification: Explicitly links Gramcianinov's regions to Reboita's RG nomenclature to establish bibliographical consistency without changing the definition source.
% --- ESTRATIFICAÇÃO GEOGRÁFICA E RESPOSTA AO FORÇAMENTO CLIMÁTICO ---
Esta estratificação geográfica é crítica para o delineamento experimental, já que estes focos respondem a regimes dinâmicos distintos. Evidência disso são as projeções de \citet{dejesus2021multimodel}, que demonstram que estas regiões apresentam sensibilidades opostas perante o forçamento climático: enquanto a região austral (ARG) experimentará um aumento na frequência de ciclones devido ao deslocamento dos oestes em direção ao polo, projeta-se uma diminuição nos focos subtropicais (SBR e LPB), confirmando que são sistemas governados por mecanismos que podem ser independentes entre si.
% De Jesus et al. (2021), Conclusions: "The projections indicate a reduction in the cyclogenesis frequency over the RG1 and RG2 regions... consistent with the poleward shift of the storm tracks. In contrast, an increase in cyclogenesis is projected for the RG3 region (southern Argentina) in the far future."

% --- INGREDIENTES DINÂMICOS E O PAPEL DA UMIDADE ---
A variabilidade da atividade ciclônica nestes focos é modulada pela estrutura da corrente de jato e seu acoplamento com as camadas baixas. \citet{vera2002cold} demonstraram que as perturbações de escala sinótica se propagam ao longo de guias de onda associadas ao jato polar e subtropical; no entanto, ao atravessar os Andes, estas perturbações experimentam mudanças significativas em sua estrutura vertical. Tal desacoplamento e posterior realinhamento entre níveis altos e baixos favorece configurações com divergência em níveis superiores, condição-chave para a intensificação ciclônica e o desenvolvimento de precipitação intensa na costa leste da América do Sul.
% Vera et al. (2002): "Cold season synoptic-scale waves over subtropical South America... The results describe the propagation of wave packets... along the subtropical and polar jets."
% --- TRANSPORTE DE UMIDADE E FORÇANTE DIABÁTICO ---
Simultaneamente, o transporte meridional de umidade desde os trópicos, mediado pelo Jato de Baixos Níveis da América do Sul (SALLJ), desempenha um papel determinante na alimentação da precipitação extrema. Como assinalam \citet{reboita2022from}, a advecção de ar quente e úmido no flanco oriental do ciclone não apenas instabiliza a atmosfera, mas atua como uma fonte diabática que, ao liberar calor latente, reforça a intensidade do sistema.
% Reboita et al. (2021): Discusses the "advection of warm and moist air" and "diabatic heating" associated with the transition of systems like Raoni.
% (Complementario: Machado 2020 sobre SALLJ también respalda esto, pero Reboita 2021 es más directo sobre la ciclogénesis en RG1).
% --- COMPLEXIDADE TERMODINÂMICA E NÃO-LINEARIDADE DOS IMPACTOS ---
No entanto, este acoplamento termodinâmico introduz uma complexidade crítica: \citet{sinclair2023relationship} advertem que a relação entre a profundidade dinâmica do sistema (vorticidade) e sua resposta hidrológica (precipitação) não é linear, já que um ciclone pode ser dinamicamente intenso, mas gerar pouca chuva se carecer deste suprimento de umidade. Este fato físico sublinha que os máximos de impacto não necessariamente coincidem com o mínimo de pressão, reforçando a necessidade de caracterizar a pegada espacial de vento e precipitação.
% Sinclair & Catto (2023): "We quantify the linear relationship between maximum vorticity and ETC-related precipitation... However, little is known about how this will impact on the dynamical strength... and whether the impact will differ for different types of ETCs." (Clave para justificar la "no linealidad").

% --- DESACOPLAMENTO ESPACIAL E MAGNITUDE DOS EXTREMOS ---
Um desafio central na caracterização dinâmica destes sistemas é a assimetria espacial de sua intensidade em superfície. A literatura tradicional tendeu a utilizar principalmente a pressão mínima central para classificar a força dos ciclones; no entanto, investigações recentes demonstram um desacoplamento sistemático entre a profundidade do vórtice e a magnitude de seus extremos associados. Estudos de \citet{bitencourt2010relating} e \citet{cardoso2022synoptic} evidenciam que os máximos de vento não se distribuem uniformemente ao redor do centro, mas concentram-se em setores específicos (como o quadrante nordeste na região subtropical) dependendo da latitude e da maturidade do sistema.
% Bitencourt (2010): "Generally, the winds are well associated with the extratropical cyclones only south of 28S... Altitude also plays an important role..." (Muestra que la relación física vorticidad-viento varía espacialmente).
% Cardoso (2022): "indicate extreme wind speed in the east and northeast of subtropical cyclogenesis region (RG1)..." (Valida la asimetría física de los máximos).
% --- DESCONEXÃO DE MÉTRICAS E EVENTOS COMPOSTOS ---
Esta desconexão sugere que a métrica de pressão é insuficiente para capturar a severidade dinâmica do evento, especialmente quando se consideram os \textit{eventos compostos}. Como mostram \citet{chen2025characteristics}, a caracterização física destes eventos emerge da co-ocorrência temporal e local de extremos de vento e precipitação, uma sobreposição que não é capturada ao analisar cada variável de forma individual.
% Chen et al. (2025), Section 2: Data and Methods % "A compound wind–precipitation extreme is defined when a local wind extreme occurs within ±6 hours of a local precipitation extreme at the same grid point."
% Chen (2025): "In NNA, 20% of local wind extremes occur +/- 6 hr of precipitation extreme... with 90% attributed to ETCs..." (Enfoca en la física de la coincidencia de eventos).
% --- FATORES CINEMÁTICOS E VELOCIDADE DE PROPAGAÇÃO ---
Adicionalmente, fatores cinemáticos como a velocidade de propagação modulam a duração e acumulação destes extremos, tal como discutem \citet{gramcianinov2023impact} ao analisar a interação do vento com a superfície oceânica.
% (Referencia a la modulación dinámica de la severidad por la velocidad de fase).
% --- CARACTERIZAÇÃO DA PEGADA ESPACIAL E EXTREMOS ---
Em consequência, basear a avaliação unicamente na trajetória do centro geométrico limita a compreensão da estrutura do fenômeno. Torna-se necessário transitar para uma caracterização da pegada espacial completa, permitindo assim quantificar com precisão a distribuição dos máximos tanto no regime climatológico médio quanto nos 10\% dos sistemas mais extremos e, adicionalmente, em uma amostra de 10\% de ciclones em torno da mediana da vorticidade máxima do ciclo de vida.

% --- PÁRRAFO 6: 
% --- EVOLUÇÃO TEMPORAL E CLASSIFICAÇÃO DINÂMICA (FASES) ---
Esta complexidade estrutural dos ciclones extratropicais não é estática; evolui drasticamente ao longo do ciclo de vida do sistema. Historicamente, a classificação destes eventos baseou-se no diagrama de espaço de fase de \citet{hart2003cyclone}, o qual discrimina sistemas segundo sua assimetria térmica e seu núcleo (frio/quente).
% Hart (2003): "An objectively defined three-dimensional cyclone phase space is proposed... classifying... cold- versus warm-core structure."
% --- LIMITAÇÕES OPERATIVAS E SISTEMAS HÍBRIDOS ---
No entanto, esta abordagem apresenta limitações operativas no Atlântico Sul, onde \citet{reboita2022from} documentaram frequentes transições rápidas e sistemas híbridos que desafiam as categorias binárias clássicas.
% Reboita (2021): "From a Shapiro-Keyser extratropical cyclone to the subtropical cyclone Raoni... The eastern coast of South America is a cyclogenetic region... of subtropical cyclones."
% --- CICLOPHASER E FASES DINÂMICAS ---
Para superar essas barreiras, esta investigação integra a recente climatologia evolutiva desenvolvida por \citet{coutodesouza2024new}, os quais, mediante o algoritmo \textit{CycloPhaser}, conseguiram segmentar objetivamente a história dos ciclones regionais em quatro fases dinâmicas discretas: incipiente, intensificação, maturidade e decaimento.
% Souza (2024): "Utilizing the minimum relative vorticity series and its derivative... the program effectively identifies distinct phases..." (Tu uso: tomas esta segmentación ya hecha para tus análisis).
% --- USO DE BASE DE DADOS CLASSIFICADA E ESTRATÉGIAS GLOBAIS ---
O uso desta base de dados classificada permite, pela primeira vez, investigar como se transforma a distribuição espacial das variáveis durante o ciclo de vida do sistema, uma estratégia alinhada com outras novas propostas globais como \textit{SyCLoPS} de \citet{han2025system} no Atlântico Norte, que enfatizam a urgência de vincular a evolução física do sistema com suas contribuições específicas de impacto.
% Han (2025): "The framework is useful to study the frequency, structure, development, wind impact... of each type of LPS." (Valida tu estrategia de vincular fase con impacto).

% --- PÁRRAFO 7: 
% --- ESTRATÉGIA ESTATÍSTICA (MEOF E REDUÇÃO DE DIMENSIONALIDADE) ---
Para caracterizar esta complexidade multidimensional, requerem-se técnicas estatísticas que transcendam a análise univariada clássica. Neste sentido, \citet{eisenstein2023identification} demonstraram que a variabilidade dos impactos extremos em ciclones pode ser reduzida a um número limitado de modos recorrentes utilizando técnicas de decomposição ortogonal.
% Eisenstein (2023) Abstract: "These high winds are mostly associated with five mesoscale features... The timing within the cyclone's life cycle... [varies]." (Valida que los ciclones complejos pueden descomponerse en "features" o modos principales).
% --- METODOLOGIA MEOF E VARIABILIDADE ACOPLADA ---
Seguindo este princípio, esta tese implementa a metodologia de Funções Ortogonais Empíricas Multivariadas (MEOF), fundamentada matematicamente por \citet{liang2018multivariate}. Diferentemente dos EOF tradicionais, esta técnica permite extrair modos de variabilidade acoplados entre campos fisicamente distintos, revelando a coerência espacial oculta onde a circulação do vento e os núcleos de precipitação covariam sistematicamente.
% Liang (2018) Abstract: "The MEOF approach allows for the extraction of coupled patterns of variability shared among the variables... constructed from nitrate, salinity, and potential temperature..." (Valida el uso de MEOF para encontrar patrones compartidos entre variables distintas).
% --- PROTÓTIPOS DE IMPACTO E SÍNTESE ESTRUTURAL ---
A aplicação deste arcabouço estatístico sobre as fases dinâmicas previamente classificadas permitirá construir "protótipos de impacto", sintetizando a estrutura típica do ciclone para cada região de gênese e estágio evolutivo, superando assim a visão estática das médias climatológicas.

% ---% ---% ---% ---% ---% ---% ---% ---% ---% ---% ---% ---% ---% ---% ---% ---% ---% ---% ---% ---
% ---% ---% ---% ---% ---% ---% ---% ---% ---% ---% ---% ---% ---% ---% ---% ---% ---% ---% ---% ---
% --- PÁRRAFO 8: 
% --- DADOS E FERRAMENTAS (ERA5 E TRACKING POR VORTICIDADE) ---
A viabilidade desta caracterização estrutural reside na recente disponibilidade do reanálise ERA5 \citep{hersbach2020era5}, cuja alta resolução espacial ($\sim$31 km) e temporal (horária) permite resolver os gradientes de pressão e os fluxos de umidade que gerações anteriores de dados suavizavam. Validações específicas para o Atlântico Sul realizadas por \citet{gramcianinov2020analysis} confirmam que o ERA5 oferece uma representação superior da intensidade dos ventos superficiais e da estrutura dos ciclones em comparação com produtos anteriores como o CFSR.
% Gramcianinov (2020): Validates ERA5 superiority in the South Atlantic.
% --- RASTREAMENTO POR VORTICIDADE E CONSENSO METODOLÓGICO ---
Para capitalizar esta precisão, este estudo utiliza a base de dados de trajetórias de \citet{coutodesouza2024new}, a qual se fundamenta no critério da vorticidade relativa em 850 hPa. Esta escolha metodológica possui um consenso transversal na literatura: estabelecida classicamente por \citet{sinclair1994objective} e consolidada por \citet{padilhareinke2024objective} como o padrão mais robusto para o Hemisfério Sul, este critério é inclusive utilizado em estudos recentes do Hemisfério Norte \citep{chen2025characteristics}.
% Couto de Souza (2024): Database provider.
% Sinclair (1994) & Padilha Reinke (2024): Theoretical justification (Classic + SH context).
% % Padilha Reinke (2024) Chen (2025) & Han (2024): Modern validation (NH context), showing global methodology consensus.
% --- VANTAGENS DO RASTREAMENTO POR VORTICIDADE E SÍNTESE TÉCNICA ---
Diferentemente dos algoritmos baseados exclusivamente na pressão mínima, o rastreamento por vorticidade permite a análise do ciclo de vida dos eventos, capturando inclusive aqueles sistemas rápidos que os algoritmos isobáricos tendem a omitir. Esta combinação de dados de alta fidelidade com uma detecção dinâmica sensível constitui a base técnica necessária para alimentar a análise multivariada (MEOF) proposta, permitindo quantificar finalmente a evolução da pegada espacial dos eventos compostos na região.
% ---% ---% ---% ---% ---% ---% ---% ---% ---% ---% ---% ---% ---% ---% ---% ---% ---% ---% ---% ---
% ---% ---% ---% ---% ---% ---% ---% ---% ---% ---% ---% ---% ---% ---% ---% ---% ---% ---% ---% ---
\section{Objetivos do Projeto de Pesquisa}

\subsection{Objetivo Geral}
Caracterizar a estrutura interna e a covariabilidade espacial dos campos de vento e precipitação nos ciclones extratropicais do Atlântico Sul ao longo de seu ciclo de vida, integrando a classificação por fases evolutivas com análise multivariada para estabelecer protótipos físicos diferenciados por região de gênese.

\subsection{Objetivos Específicos}
\begin{enumerate}
    \item Consolidar uma base de dados lagrangiana que associe os campos de superfície de alta resolução (ERA5) com as trajetórias e fases dinâmicas (incipiente, intensificação, maturidade, decaimento) identificadas para as três regiões de ciclogênese (SBR, LPB, ARG).
    \item Determinar a assimetria espacial dos extremos mediante a estimativa de densidade por núcleos (KDE), quantificando a probabilidade de ocorrência de máximos de vento e precipitação em relação ao centro do ciclone para identificar o desacoplamento físico entre a pressão mínima e a severidade em superfície.
    \item Identificar os modos de variabilidade acoplada aplicando Funções Ortogonais Empíricas Multivariadas (MEOF) aos campos conjuntos de vento e precipitação, para revelar as estruturas coerentes de mesoescala que governam os eventos compostos em cada fase evolutiva.
    \item Construir modelos conceituais (protótipos) para cada região de gênese, sintetizando os padrões dominantes obtidos (MEOF) e as distribuições de probabilidade de extremos (KDE) em uma representação composta que descreva a evolução em suas fases de vida.
\end{enumerate}